% !TeX root = TIPE_vangi.tex
% !TeX spellcheck = fr

\section{Simulations informatiques}

\subsection{Représentation des différentes données}

Tous les ensembles sont représenté par des listes de chaînes de caractères.
Une préordre $\theta$ sont représenté pour $\varXi_\theta(\C)$ (qui est une liste de liste de chaînes de caractères) dans la classe \texttt{Preference}.

La classe abstraite \texttt{Vote} décrit tous les systèmes de vote possibles avec les méthodes \texttt{\_\_init\_\_, voter, vainqueur, resultat, copier, \_\_repr\_\_}. Les systèmes de vote concrets en sont des enfants.

\bigbreak
\begin{figure}[H]
	\centering
	\begin{forest}
		for tree={s sep=5mm, l=2cm}
		[\texttt{Vote}
			[\texttt{Pondere}
				[\texttt{Uninominal}]
				[\texttt{Borda}]
				[\texttt{Approbation}]
			]
			[\texttt{Condorcet}
				[\texttt{Black}]
				[\texttt{Copeland}]
				[\texttt{Schulze}]
				[\texttt{GagnantCondorcet}]
			]
		]
	\end{forest}
\end{figure}
\bigbreak

Il faut calculer différents ensembles ($\C$, $\pO$, $\mathfrak{O}$), l’ensemble des applications ${\D}^\V$ et leurs systèmes de représentants.
Cependant le calcul de leur cardinal est souvent bien plus simple que le calcul explicite des ensembles.
C'est pourquoi ils sont décrits par la classe abstraite \texttt{Ensemble} que possède les méthodes \texttt{\_\_init\_\_, card = \_\_len\_\_, \_\_call\_\_}. Les classes fils implémentent les méthodes \texttt{\_calcul\_card} et \texttt{\_calcul\_ensemble} qui sont appelées, si besoin, puis mémoïsé par la classe mère.

Le calcul de la probabilité se base sur \hyperref[individuel d'invariance par impartialité]{le théorème d'invariance par impartialité}.
La somme sur $\mathcal{S}$ permet de paralléliser les calculs.

Le code Python final fait 1500 lignes.

\subsection{Graphes obtenus}
\def \sgh {0.6} % scale graphe

Nous allons étudier cinq systèmes de vote \textit{le scrutin uninominal, le méthode Borda, la méthode Black, le méthode de Copeland et la méthode de Schulze}.
Le nombre d'électeurs est en abscisse et la probabilité en ordonnée.
Les courbes bleus correspondent à deux choix, celles oranges à trois choix, celles vertes à quatre choix et celles rouges à cinq choix.
Les courbes bleus sont identiques puisque, pour deux choix, tous les systèmes de vote étudiés coïncident avec le référendum.
Sauf mention contraire, $\D = \mathfrak{O}$.
\newline

\begin{figure}[H]
	\underline{Critère d'honnêteté}
	\begin{adjustwidth}{-1cm}{-1cm}
		\centering
		\begin{subfigure}{1\linewidth}
			\includegraphics[scale=\sgh]{honnetete/graphe_Uninominal_ordres_3_30.png}
			\hfill
			\includegraphics[scale=\sgh]{honnetete/graphe_Borda_ordres_3_30.png}
		\end{subfigure}
	\end{adjustwidth}
\end{figure}

\begin{figure}[H]
	\begin{adjustwidth}{-1cm}{-1cm}
		\centering
		\begin{subfigure}{1\linewidth}
			\includegraphics[scale=\sgh]{honnetete/graphe_Black_ordres_3_30.png}
			\hfill
			\includegraphics[scale=\sgh]{honnetete/graphe_Copeland_ordres_3_30.png}
		\end{subfigure}
	\end{adjustwidth}
\end{figure}
\vspace{-1cm}
\begin{figure}[H]
	\begin{adjustwidth}{-1cm}{-1cm}
		\begin{subfigure}{1\linewidth}
			\includegraphics[scale=\sgh]{honnetete/graphe_Schulze_ordres_3_25.png}
		\end{subfigure}
	\end{adjustwidth}
\end{figure}

\begin{figure}[H]
	\underline{Critère de Condorcet} \\
	Pour les systèmes de Condorcet, la probabilité vaut bien sûr 1.
	\begin{adjustwidth}{-1cm}{-1cm}
		\centering
		\begin{subfigure}{1\linewidth}
			\includegraphics[scale=\sgh]{critereCondorcet/graphe_Uninominal_ordres_3_30.png}
			\hfill
			\includegraphics[scale=\sgh]{critereCondorcet/graphe_Borda_ordres_3_30.png}
		\end{subfigure}
	\end{adjustwidth}
\end{figure}

\begin{figure}[H]
	\underline{Critère d'unicité du vainqueur}
	\begin{adjustwidth}{-1cm}{-1cm}
		\centering
		\begin{subfigure}{1\linewidth}
			\includegraphics[scale=\sgh]{unique_vainqueur/graphe_Uninominal_ordres_3_30.png}
			\hfill
			\includegraphics[scale=\sgh]{unique_vainqueur/graphe_Borda_ordres_3_30.png}
		\end{subfigure}
	\end{adjustwidth}
\end{figure}
\vspace{-1cm}
\begin{figure}[H]
	\begin{adjustwidth}{-1cm}{-1cm}
		\centering
		\begin{subfigure}{1\linewidth}
			\includegraphics[scale=\sgh]{unique_vainqueur/graphe_Black_ordres_3_30.png}
			\hfill
			\includegraphics[scale=\sgh]{unique_vainqueur/graphe_Copeland_ordres_3_30.png}
		\end{subfigure}
	\end{adjustwidth}
\end{figure}
\vspace{-1cm}
\begin{figure}[H]
	\begin{adjustwidth}{-1cm}{-1cm}
		\begin{subfigure}{1\linewidth}
			\includegraphics[scale=\sgh]{unique_vainqueur/graphe_Schulze_ordres_3_20.png}
		\end{subfigure}
	\end{adjustwidth}
\end{figure}

\begin{figure}[H]
	\underline{Critère de particpation} \\
	La probabilité de ce critère vaut 1 pour les votes pondérés d'après \hyperref[Vote pondéré croissant]{le théorème des votes pondérés croissants}.
	\begin{adjustwidth}{-1cm}{-1cm}
		\centering
		\begin{subfigure}{1\linewidth}
			\includegraphics[scale=\sgh]{participation/graphe_Black_ordres_3_30.png}
			\hfill
			\includegraphics[scale=\sgh]{participation/graphe_Copeland_ordres_3_30.png}
		\end{subfigure}
	\end{adjustwidth}
\end{figure}
\vspace{-1cm}
\begin{figure}[H]
	\begin{adjustwidth}{-1cm}{-1cm}
		\begin{subfigure}{1\linewidth}
			\includegraphics[scale=\sgh]{participation/graphe_Schulze_ordres_3_20.png}
		\end{subfigure}
	\end{adjustwidth}
\end{figure}

\begin{figure}[H]
	\underline{Avec plus de trois candidats} \\
	\begin{adjustwidth}{-1cm}{-1cm}
		\centering
		\begin{subfigure}{1\linewidth}
			\includegraphics[scale=\sgh]{honnetete/graphe_Uninominal_ordres_4_6.png}
			\hfill
			\includegraphics[scale=\sgh]{honnetete/graphe_Borda_ordres_4_6.png}
		\end{subfigure}
	\end{adjustwidth}
\end{figure}
\vspace{-1cm}
\begin{figure}[H]
	\begin{adjustwidth}{-1cm}{-1cm}
		\centering
		\begin{subfigure}{1\linewidth}
			\includegraphics[scale=\sgh]{honnetete/graphe_Black_ordres_4_6.png}
			\hfill
			\includegraphics[scale=\sgh]{honnetete/graphe_Copeland_ordres_4_6.png}
		\end{subfigure}
	\end{adjustwidth}
\end{figure}
\vspace{-1cm}
\begin{figure}[H]
	\begin{adjustwidth}{-1cm}{-1cm}
		\begin{subfigure}{1\linewidth}
			\includegraphics[scale=\sgh]{honnetete/graphe_Schulze_ordres_4_5.png}
		\end{subfigure}
	\end{adjustwidth}
\end{figure}

\begin{figure}[H]
	\begin{adjustwidth}{-1cm}{-1cm}
		\centering
		\begin{subfigure}{1\linewidth}
			\includegraphics[scale=\sgh]{unique_vainqueur/graphe_Uninominal_ordres_4_7.png}
			\hfill
			\includegraphics[scale=\sgh]{unique_vainqueur/graphe_Borda_ordres_4_7.png}
		\end{subfigure}
	\end{adjustwidth}
\end{figure}
\begin{figure}[H]
	\begin{adjustwidth}{-1cm}{-1cm}
		\centering
		\begin{subfigure}{1\linewidth}
			\includegraphics[scale=\sgh]{unique_vainqueur/graphe_Black_ordres_4_6.png}
			\hfill
			\includegraphics[scale=\sgh]{unique_vainqueur/graphe_Copeland_ordres_4_6.png}
		\end{subfigure}
	\end{adjustwidth}
\end{figure}

\begin{figure}[H]
	\underline{Avec de relations de préordre} \\
	Maintenant, $\D = \pO$.
	\begin{adjustwidth}{-1cm}{-1cm}
		\centering
		\begin{subfigure}{1\linewidth}
			\includegraphics[scale=\sgh]{honnetete/graphe_Borda_préordres_3_10.png}
			\hfill
			\includegraphics[scale=\sgh]{honnetete/graphe_Black_préordres_3_9.png}
		\end{subfigure}
	\end{adjustwidth}
\end{figure}
\vspace{-1cm}
\begin{figure}[H]
	\begin{adjustwidth}{-1cm}{-1cm}
		\centering
		\begin{subfigure}{1\linewidth}
			\includegraphics[scale=\sgh]{honnetete/graphe_Copeland_préordres_3_9.png}
			\hfill
			\includegraphics[scale=\sgh]{honnetete/graphe_Schulze_préordres_3_8.png}
		\end{subfigure}
	\end{adjustwidth}
\end{figure}

\begin{figure}[H]
	\begin{adjustwidth}{-1cm}{-1cm}
		\centering
		\begin{subfigure}{1\linewidth}
			\includegraphics[scale=\sgh]{unique_vainqueur/graphe_Borda_préordres_3_10.png}
			\hfill
			\includegraphics[scale=\sgh]{unique_vainqueur/graphe_Black_préordres_3_9.png}
		\end{subfigure}
	\end{adjustwidth}
\end{figure}
\vspace{-1cm}
\begin{figure}[H]
	\begin{adjustwidth}{-1cm}{-1cm}
		\centering
		\begin{subfigure}{1\linewidth}
			\includegraphics[scale=\sgh]{unique_vainqueur/graphe_Copeland_préordres_3_9.png}
			\hfill
			\includegraphics[scale=\sgh]{unique_vainqueur/graphe_Schulze_préordres_3_8.png}
		\end{subfigure}
	\end{adjustwidth}
\end{figure}

\begin{figure}[H]
	\begin{adjustwidth}{-1cm}{-1cm}
		\centering
		\begin{subfigure}{1\linewidth}
			\includegraphics[scale=\sgh]{participation/graphe_Black_préordres_3_9.png}
			\hfill
			\includegraphics[scale=\sgh]{participation/graphe_Copeland_préordres_3_9.png}
		\end{subfigure}
	\end{adjustwidth}
\end{figure}
\vspace{-1cm}
\begin{figure}[H]
	\begin{adjustwidth}{-1cm}{-1cm}
		\begin{subfigure}{1\linewidth}
			\includegraphics[scale=\sgh]{participation/graphe_Schulze_préordres_3_10.png}
		\end{subfigure}
	\end{adjustwidth}
\end{figure}

\begin{figure}[H]
	\underline{Vote par approbation} \\
	Ici, $\D = \{ \theta \in \pO \;|\; \left| \varXi_\theta(\C) \right| \leqslant 2 \}$.
	Le vote par approbation respecte le critère de participation et celui de Condorcet.
	\begin{adjustwidth}{-1cm}{-1cm}
		\centering
		\begin{subfigure}{1\linewidth}
			\includegraphics[scale=\sgh]{honnetete/graphe_Approbation_préordresduals_3_20.png}
			\hfill
			\includegraphics[scale=\sgh]{unique_vainqueur/graphe_Approbation_préordresduals_3_25.png}
		\end{subfigure}
	\end{adjustwidth}
\end{figure}
\vspace{-1cm}
\begin{figure}[H]
	\begin{adjustwidth}{-1cm}{-1cm}
		\centering
		\begin{subfigure}{1\linewidth}
			\includegraphics[scale=\sgh]{honnetete/graphe_Approbation_préordresduals_5_5.png}
			\hfill
			\includegraphics[scale=\sgh]{unique_vainqueur/graphe_Approbation_préordresduals_5_6.png}
		\end{subfigure}
	\end{adjustwidth}
\end{figure}
\vspace{-1cm}
\begin{figure}[H]
	\noindent Dans le même domaine de vote, pour deux autres systèmes de vote :
	\begin{adjustwidth}{-1cm}{-1cm}
		\centering
		\begin{subfigure}{1\linewidth}
			\includegraphics[scale=\sgh]{honnetete/graphe_Copeland_préordresduals_3_15.png}
			\hfill
			\includegraphics[scale=\sgh]{honnetete/graphe_Schulze_préordresduals_3_10.png}
		\end{subfigure}
	\end{adjustwidth}
\end{figure}